% ============================================================
% Préambule — PT Foundations (article FR, 40 p)
% Adapté de PT_MONOGRAPHY/preamble.tex
% ============================================================

\documentclass[11pt,a4paper]{article}

% --- Encodage et langue ---
\usepackage[utf8]{inputenc}
\usepackage[T1]{fontenc}
\usepackage[french]{babel}
\usepackage{csquotes}

% --- Mathématiques ---
\usepackage{amsmath,amssymb,amsthm,mathtools}
\usepackage{bm}

% --- Mise en page ---
\usepackage[margin=2.5cm]{geometry}
\usepackage{microtype}
\usepackage{enumitem}

% --- Tableaux et figures ---
\usepackage{booktabs}
\usepackage{graphicx}
\usepackage{tikz}
\usetikzlibrary{arrows.meta,positioning,fit,calc}

% --- Bibliographie ---
\usepackage[backend=biber,style=numeric,sorting=none]{biblatex}
\addbibresource{references.bib}

% --- Hyperliens ---
\usepackage[colorlinks=true,linkcolor=blue!60!black,citecolor=green!50!black,urlcolor=blue!50!black]{hyperref}
\usepackage{cleveref}

% --- Environnements de théorème ---
\theoremstyle{plain}
\newtheorem{theoreme}{Théorème}[section]
\newtheorem{proposition}[theoreme]{Proposition}
\newtheorem{lemme}[theoreme]{Lemme}
\newtheorem{corollaire}[theoreme]{Corollaire}

\theoremstyle{definition}
\newtheorem{definition}[theoreme]{Définition}
\newtheorem{convention}[theoreme]{Convention}

\theoremstyle{remark}
\newtheorem{remarque}[theoreme]{Remarque}
\newtheorem{exemple}[theoreme]{Exemple}
\newtheorem*{conjecture*}{Conjecture}
\newtheorem{conjecture}[theoreme]{Conjecture}

% --- Macros PT canoniques ---
% Secteurs (convention 2026 : q_+ / q_-, jamais q_stat/q_therm)
\newcommand{\qplus}{q_{+}}
\newcommand{\qminus}{q_{-}}

% Constantes structurelles
\newcommand{\mustar}{\mu^{*}}
\newcommand{\sper}{s}

% Dimensions anomales
\newcommand{\gammap}[1]{\gamma_{#1}}
\newcommand{\gammathree}{\gamma_{3}}
\newcommand{\gammafive}{\gamma_{5}}
\newcommand{\gammaseven}{\gamma_{7}}

% Phases cycliques
\newcommand{\thetap}[1]{\theta_{#1}}
\newcommand{\sinp}[1]{\sin^{2}\theta_{#1}}

% Status tags (10/10 [THM])
\newcommand{\THM}{\textsc{[Thm]}}

% --- Théorèmes T-series (nommés selon NOMENCLATURE_MAP.md) ---
\newcommand{\Tzero}{\textbf{T0}}
\newcommand{\Tone}{\textbf{T1}}
\newcommand{\Ttwo}{\textbf{T2}}
\newcommand{\Tthree}{\textbf{T3}}
\newcommand{\Tfour}{\textbf{T4}}
\newcommand{\Tfive}{\textbf{T5}}
\newcommand{\Tsix}{\textbf{T6}}

% --- Axiomes de pont BA-series (théorèmes prouvés) ---
\newcommand{\BAzero}{\textbf{BA0}}
\newcommand{\BAone}{\textbf{BA1}}
\newcommand{\BAtwo}{\textbf{BA2}}
\newcommand{\BAthree}{\textbf{BA3}}
\newcommand{\BAfour}{\textbf{BA4}}
\newcommand{\BAfive}{\textbf{BA5}}

% --- Lemme L0 (entropie maximale) ---
\newcommand{\Lzero}{\textbf{L0}}

% --- Macros sp\'ecifiques Article 2 ---
\newcommand{\Cper}{\mathcal{C}_{\rm per}}
\newcommand{\genus}{g}
\newcommand{\gPT}{g_{\rm PT}}
\newcommand{\Ric}{\mathrm{Ric}}
\newcommand{\Hess}{\mathrm{Hess}}
\newcommand{\Rres}{R}
\newcommand{\HBK}{H_{\rm BK}}
\newcommand{\lambdaH}{\lambda_H}
\newcommand{\DeltaMaslov}{\Delta_N^{\rm Maslov}}

% --- M\'etadonn\'ees article ---
\title{La Th\'eorie de la Persistance : lectures g\'eom\'etrique et spectrale}
\author{Yan Senez}
\date{\today}


\begin{document}

\maketitle

\begin{abstract}
\noindent
La Théorie de la Persistance (PT), exposée dans son article compagnon~\cite{PT_Article1_Foundations}, dérive d'un unique axiome arithmétique~$\sper = 1/2$ une structure de cascade sur les nombres premiers actifs~$\{3, 5, 7\}$ et un point fixe arithmétique~$\mustar = 15$. Cet article second explore trois lectures complémentaires de cette structure : (i) une lecture \emph{algébro-géométrique} qui exhibe la cascade canonique comme l'unique courbe algébrique dans la classe de Kontsevich--Norbury satisfaisant trois conditions d'holonomie, d'activation et d'auto-cohérence, et calcule son genre $\genus = \mustar - 1 = 14$ ; (ii) une lecture \emph{différentielle} qui établit que la métrique de Fisher--Bianchi induite par la cascade satisfait l'équation du soliton de Ricci $\Ric + \Hess(f) = \lambda(\mu)\, \gPT$ avec $\lambda(\mu) = -1/\mu^{4} + O(1/\mu^{5})$, coefficient~$-1$ exact ; (iii) une lecture \emph{spectrale} qui établit l'identité triple $1/8 = \sper^{2}/2 = \lambdaH = \DeltaMaslov$ reliant le couplage de Higgs et la phase de Maslov des zéros de Riemann à l'unique entrée formelle de la PT, ainsi qu'une identité de trace exacte vérifiée à $10^{-26}$. Le programme spectral-arithmétique plus ambitieux dont ces identités sont les premières pierres (conjecture RH-PT) est esquissé en perspective.
\end{abstract}

\tableofcontents

\section{Introduction et rappels}
\label{sec:intro2}

\subsection{Pourquoi des lectures géométrique et spectrale ?}
\label{sec:intro2:why}

L'article compagnon~\cite{PT_Article1_Foundations} expose les fondations purement arithmétiques de la Théorie de la Persistance (PT) : un unique axiome $\sper = 1/2$, sept théorèmes structurants $\Tzero$--$\Tsix$, trois axiomes de pont prouvés $\BAthree$--$\BAfive$, un point fixe $\mustar = 15$ et une cascade canonique sur les premiers actifs $\{3,5,7\}$. Cette structure produit, sans paramètre ajustable, des observables physiques telles que la constante de structure fine $\alpha_{\rm EM}$ à $0{,}004$~ppb.

Une fois ce socle arithmétique posé, trois questions naturelles se posent.
\begin{enumerate}[nosep,leftmargin=*]
\item \emph{La cascade est-elle un objet algébro-géométrique~?} La séquence $\gammathree \to \gammathree\gammafive \to \gammathree\gammafive\gammaseven$ admet-elle une réalisation comme courbe algébrique dont l'invariant topologique encode $\mustar$~?
\item \emph{La cascade est-elle un objet géométrique différentiel~?} La structure de l'information induit-elle une métrique riemannienne canonique sur l'espace des paramètres, et cette métrique vérifie-t-elle une équation de Ricci~?
\item \emph{La cascade est-elle un objet spectral~?} L'axiome $\sper = 1/2$, qui fixe la valeur propre dominante de la matrice de transfert primorielle (\Ttwo), produit-il des signatures spectrales identifiables dans la phénoménologie physique ou dans la structure des zéros de Riemann~?
\end{enumerate}
Cet article répond positivement aux trois questions, en présentant un résultat principal par lecture.

\subsection{Rappels minimaux}
\label{sec:intro2:rappels}

Le lecteur qui n'aurait pas lu~\cite{PT_Article1_Foundations} trouvera ci-dessous les éléments strictement nécessaires à la suite. Pour la justification, on se reportera à l'article compagnon.

\paragraph{Axiome et point fixe.} L'unique entrée formelle de la PT est $\sper = 1/2$. Le point fixe arithmétique du crible primoriel, déterminé par le théorème~\Tfive, vaut
\[
\mustar \;=\; 3 + 5 + 7 \;=\; 15.
\]
L'ensemble des nombres premiers se partage canoniquement en quatre classes :
\[
\underbrace{\{3, 5, 7\}}_{\text{actifs}}, \quad
\underbrace{\{11, 13\}}_{\text{échos}}, \quad
\underbrace{\{17, 19, 23\}}_{\text{super-échos}}, \quad
\underbrace{\{p \geq 29\}}_{\text{inactifs}},
\]
plus le canal de parité $\{2\}$.

\paragraph{Cascade canonique.} La cascade canonique de la PT (\cite[\S{}5]{PT_Article1_Foundations}) est la séquence ordonnée
\[
\mathcal C_0 \;\xrightarrow{\;\gammathree\;}\; \mathcal C_3
\;\xrightarrow{\;\gammafive\;}\; \mathcal C_{3,5}
\;\xrightarrow{\;\gammaseven\;}\; \mathcal C_{3,5,7},
\]
où $\gammap{p}(\mu) = 1 - 2/(p(1 - q^{p}))$, $q = 1 - 2/\mu$, sont les \emph{dimensions anomales}. À $\mustar = 15$, on a $\gammathree = 0{,}808$, $\gammafive = 0{,}696$, $\gammaseven = 0{,}595$.

\paragraph{Holonomie et secteurs.} La formule d'holonomie ($\BAthree$) donne, pour chaque premier actif,
\[
\sinp{p}(q) \;=\; \delta_p(q)\,\bigl(2 - \delta_p(q)\bigr), \qquad \delta_p(q) = \dfrac{1 - q^{p}}{p}.
\]
Les deux paramétrisations canoniques de la loi géométrique au point fixe sont $\qplus = 13/15$ (branche-sommet, rationnelle, max-entropie) et $\qminus = e^{-1/15}$ (branche-arête, transcendante, Gibbs), reliées par une involution $D$.

\paragraph{Produit de Pontryagin.} La fonctionnelle de couplage du crible, au point fixe $\mustar$, prend la forme produit ($\BAfive$),
\[
\alpha_{\rm sieve} \;=\; \prod_{p \in \{3,5,7\}} \sinp{p}(\qplus).
\]

\subsection{Plan, statut épistémique}
\label{sec:intro2:plan}

La suite est organisée en cinq sections complémentaires.
\begin{itemize}[nosep,leftmargin=*]
\item La section~\ref{sec:courbe} établit le \emph{théorème de la courbe de persistance} : la cascade canonique admet une unique courbe algébrique dans la classe de Kontsevich--Norbury, satisfaisant trois conditions d'holonomie, de seuil et d'auto-cohérence. Son genre vaut $\genus = \mustar - 1 = 14$.
\item La section~\ref{sec:soliton} établit le \emph{théorème du quasi-soliton de Ricci} : la métrique de Fisher--Bianchi $\gPT$ induite par la cascade vérifie $\Ric(\gPT) + \Hess(f) = \lambda(\mu)\, \gPT$ avec $\lambda(\mu) = -1/\mu^{4} + O(1/\mu^{5})$, coefficient $-1$ exact.
\item La section~\ref{sec:spectrales} démontre l'\emph{identité spectrale K2} : $1/8 = \sper^{2}/2 = \lambdaH = \DeltaMaslov$, identifiant la phase de Maslov des zéros de Riemann avec le couplage de Higgs et l'axiome de la PT, et l'\emph{identité de trace A6} vérifiée numériquement à $10^{-26}$.
\item La section~\ref{sec:programme} esquisse le programme spectral-arithmétique plus ambitieux (\emph{conjecture RH-PT}) dont les identités K2 et A6 sont les premières pierres.
\item La section~\ref{sec:conclusion2} synthétise les trois lectures et trace les perspectives.
\end{itemize}

\textbf{Statut épistémique.} Les résultats de cet article se répartissent en trois statuts distincts qu'il convient de ne pas confondre.
\begin{description}[nosep,leftmargin=*]
\item[Théorèmes prouvés inconditionnellement] : courbe de persistance et calcul du genre (\ref{sec:courbe}, d'après \cite{PT_GeoFlow_PersistenceCurve}), quasi-soliton de Ricci asymptotique (\ref{sec:soliton}, d'après \cite{PT_GeoFlow_RicciSoliton}).
\item[Identités prouvées] : K2 et A6 sont des identités exactes (la première structurelle, vérifiée à toutes les précisions disponibles ; la seconde algébrique en $\Re(s) > 1$, vérifiée numériquement à $10^{-26}$). Leur prolongement à la ligne critique reste ouvert.
\item[Programme conjectural] : la conjecture RH-PT, qui prolongerait A6 et K2 à $\Re(s) = 1/2$, est un programme actif, non un théorème.
\end{description}
Cette différenciation est maintenue tout au long de l'article.

\section{Lecture algébro-géométrique : la courbe de persistance}
\label{sec:courbe}

La cascade canonique de la PT (\ref{sec:intro2:rappels}) est une séquence purement combinatoire de raffinements arithmétiques. Cette section montre qu'elle admet une réalisation algébro-géométrique unique : il existe une courbe algébrique dans la classe de Kontsevich--Norbury qui code exactement la cascade des premiers actifs, et son invariant topologique fondamental --- le genre --- s'exprime par la formule remarquable $\genus = \mustar - 1 = 14$. Tous les résultats de cette section sont rigoureusement prouvés dans~\cite{PT_GeoFlow_PersistenceCurve} ; nous en donnons les énoncés et les idées de preuve.

\subsection{Préliminaires : courbes spectrales de Kontsevich--Norbury}
\label{sec:courbe:KN}

Une \emph{courbe spectrale} au sens de la récursion topologique d'Eynard--Orantin~\cite{EynardOrantin2007} est la donnée d'un quadruplet $(\Sigma, x, y, B)$ où $\Sigma$ est une surface de Riemann, $x$ et $y$ deux fonctions méromorphes sur $\Sigma$, et $B$ une forme bilinéaire (le \emph{noyau de Bergman}). La classe de Kontsevich--Norbury~\cite{Kontsevich1992,Norbury2017}, notée $\mathcal{KN}$, est la sous-classe des courbes spectrales dont la forme univariée $\omega_{0,1} = y\, dx$ détermine, via la récursion topologique, l'ensemble des invariants de Weil--Petersson de l'espace des modules $\overline{\mathcal M}_{g,n}$ des surfaces de Riemann pointées. Le cas archétypique est la courbe d'Airy de Mirzakhani~\cite{Mirzakhani2007}, dont l'univariée fait apparaître la constante $2\pi^{2}/3$ caractéristique des volumes de Weil--Petersson.

\subsection{Trois conditions d'auto-cohérence}
\label{sec:courbe:conditions}

À une cascade $S = \{p_{(1)}, \dots, p_{(k)}\} \subseteq \mathbb N^{*}$ avec $\mustar(S) := \sum_{p \in S} p > 2$ on associe naturellement une famille de courbes candidates dans $\mathcal{KN}$, paramétrées par une variable complexe $z$ via
\begin{equation}
\mu_S(z) \;=\; \dfrac{\mustar(S)}{1 - z^{2}},
\qquad
q_S(z) \;=\; 1 - \dfrac{2}{\mu_S(z)},
\label{eq:param_curve}
\end{equation}
et la fonction $y$ définie par le produit cascade
\begin{equation}
y_S(z) \;=\; \dfrac{z}{2\pi} \prod_{n \in S} \sin\thetap{n}\bigl(\mu_S(z)\bigr),
\qquad
\omega_{0,1}^{\mathcal C_S} \;=\; y_S(z)\,z\,dz.
\label{eq:y_S}
\end{equation}
La paramétrisation~\eqref{eq:param_curve} est rationnelle paire en $z$, à pôles simples en $z = \pm 1$, vérifiant $\mu_S(0) = \mustar(S)$. Le point $z = 0$ est par conséquent le \emph{point d'auto-cohérence} : la valeur du paramètre $\mu$ y est précisément le point fixe arithmétique.

Pour qu'une telle courbe code authentiquement la cascade PT, trois conditions structurelles doivent être imposées simultanément.

\begin{definition}[Conditions ACA$_\mathcal N$ de cascade]
\label{def:ABC}
Soit $\mathcal N \subseteq \mathbb N^{*}$ un crible (typiquement $\mathcal N = \{$nombres premiers$\}$, mais d'autres choix sont admissibles), et soit $(S, \mu)$ avec $S \subseteq \mathcal N$ fini et $\mu : \mathbb C \to \widehat{\mathbb C}$ rationnelle paire.
\begin{enumerate}[nosep,label=\textbf{(\Alph*)}]
\item \emph{Holonomie polynomiale.} Pour chaque $n \in S$, la phase cyclique vérifie $\sinp{n} = \delta_n(q)(2 - \delta_n(q))$ avec $\delta_n(q) = (1 - q^{n})/n$ (axiome de pont $\BAthree$ de la PT).
\item \emph{Seuil d'activation.} Le sous-ensemble $S$ coïncide avec l'ensemble des éléments de $\mathcal N$ pour lesquels $\gammap{n}(\mustar) > \sper = 1/2$ (axiome $\BAfour$).
\item \emph{Auto-cohérence.} Le paramètre vérifie $\mu(0) = \mustar(S) := \sum_{n \in S} n$ (théorème~$\Tfive$).
\end{enumerate}
\end{definition}

Sur le crible des nombres premiers $\mathcal N = \mathcal P$, les conditions (A), (B), (C) sont précisément la traduction algébrique des théorèmes et axiomes de pont d'Article~1. Le théorème principal de cette section affirme que ces trois conditions sélectionnent un \emph{unique} sous-ensemble $S$ et que la courbe associée est algébrique.

\subsection{Théorème d'unicité et classification}
\label{sec:courbe:classification}

\begin{theoreme}[Classification universelle, version simplifiée]
\label{thm:courbe_classification}
Soit $\mathcal P$ l'ensemble des nombres premiers. Parmi les sous-ensembles $S \subseteq \mathcal P$ vérifiant les conditions ACA$_{\mathcal P}$ (\ref{def:ABC}), il existe \emph{exactement un} $S$ avec $|S| \geq 2$ et $\mustar(S) \leq 200$~: c'est
\[
S \;=\; \{3, 5, 7\}, \qquad \mustar(S) \;=\; 15.
\]
\end{theoreme}

\textbf{Idée de preuve.} La condition (B), $\gammap{n}(\mustar) > 1/2$, est très contraignante : combinée à la monotonie décroissante de $\gammap{n}(\mu)$ avec $n$ démontrée dans l'article compagnon~\cite{companion_M2}, elle force $S$ à être contenu dans un segment initial des premiers. La condition (C), $\mustar(S) = \sum_{n \in S} n$, équilibre alors l'effet : un $S$ trop petit ne sature pas la somme, un $S$ trop grand viole le seuil. La recherche numérique exhaustive sur tous les $2406$ sous-ensembles compatibles jusqu'à $\mustar \leq 200$ identifie $S = \{3, 5, 7\}$ comme unique solution avec $|S| \geq 2$~\cite[\S 5]{PT_GeoFlow_PersistenceCurve}. La preuve analytique en complète l'argument pour $\mustar > 200$ via les bornes du théorème des nombres premiers.

\subsection{Forme algébrique explicite}
\label{sec:courbe:forme}

\begin{theoreme}[Algébricité de la courbe de persistance]
\label{thm:algebraicite_PT}
La courbe $\mathcal C_S$ définie par~\eqref{eq:param_curve}--\eqref{eq:y_S} avec $S = \{3, 5, 7\}$ est algébrique : la fonction $y^{2}$ est une fraction rationnelle de $x = z^{2}/2$. Explicitement, en posant $q(x) = (\mustar - 2 + 4x)/\mustar$,
\begin{equation}
Y^{2} \;=\; \dfrac{2x}{4\pi^{2}} \prod_{n \in S} \dfrac{(1 - q(x)^{n})\,(2n - 1 + q(x)^{n})}{n^{2}} \;=\; \dfrac{N_S(x)}{2\pi^{2}\,D_S},
\label{eq:Y_squared}
\end{equation}
avec $D_S \in \mathbb N$ et $N_S \in \mathbb Z[x]$ de degré $2 \mustar + 1$. Le polynôme $N_S$ admet $x$ en facteur (multiplicité~$1$, du préfacteur $z^{2} = 2x$) et $(2x - 1)$ en facteur (multiplicité $|S|$, annulation simultanée des $|S|$ facteurs $(1 - q^{n})$ pour $q = 1$).
\end{theoreme}

\textbf{Idée de preuve.} La paramétrisation $\mu_S(z)$ étant rationnelle paire, $q_S(z) = 1 - 2/\mu_S(z)$ l'est aussi. Chaque $\sin^{2}\thetap{n}(\mu_S(z)) = \delta_n(2 - \delta_n)$ avec $\delta_n = (1 - q^{n})/n$ est rationnel en $q^{n}$, donc en $z^{2}$, donc en $x = z^{2}/2$. Le produit~\eqref{eq:Y_squared} est donc rationnel en $x$. Pour $S = \{3, 5, 7\}$, le degré explicite de $N_S$ vaut $2 \cdot 15 + 1 = 31$ ; le calcul symbolique exact (Python/SymPy, scripts dans~\cite{PT_GeoFlow_PersistenceCurve}) confirme la structure factorielle annoncée.

\subsection{Calcul du genre et conséquence topologique}
\label{sec:courbe:genre}

\begin{theoreme}[Formule universelle du genre]
\label{thm:genre}
Soit $S \subseteq \mathcal P$ une solution des conditions ACA$_{\mathcal P}$ vérifiant les hypothèses du \ref{thm:courbe_classification}. Alors le genre de la courbe algébrique sous-jacente à $\mathcal C_S$ est
\[
\boxed{\;\genus(\mathcal C_S) \;=\; \mustar(S) - \lfloor |S|/2 \rfloor.\;}
\]
\end{theoreme}

\textbf{Conséquence pour PT.} En appliquant la formule à l'unique solution $S = \{3, 5, 7\}$ de \ref{thm:courbe_classification},
\begin{equation}
\boxed{\;\genus(\Cper) \;=\; \mustar - \lfloor 3/2 \rfloor \;=\; 15 - 1 \;=\; 14.\;}
\label{eq:genus_14}
\end{equation}

\textbf{Idée de preuve.} On part de la forme explicite~\eqref{eq:Y_squared}, on extrait du polynôme $N_S$ son facteur squarefree $N_S^{\rm sqf}(x) = N_S(x)/(2x - 1)^{|S| - 1}$. Le degré de $N_S^{\rm sqf}$ est $2\mustar + 1 - (|S| - 1) = 2\mustar - |S| + 2$. La courbe hyperelliptique $Y^{2} = N_S^{\rm sqf}(x)$ a alors genre $\lfloor (\deg N_S^{\rm sqf} - 1)/2 \rfloor = \lfloor (2\mustar - |S| + 1)/2 \rfloor = \mustar - \lfloor |S|/2 \rfloor$. Pour $S = \{3, 5, 7\}$, calcul direct : $\genus = 15 - 1 = 14$. Le calcul exhaustif des discriminants des polynômes squarefree est dans~\cite[\S 7]{PT_GeoFlow_PersistenceCurve}.

\subsection{Interprétation : \texorpdfstring{$\mustar$}{μ*} comme invariant topologique}
\label{sec:courbe:interpretation}

L'identité~\eqref{eq:genus_14} transforme le point fixe arithmétique $\mustar = 15$ d'un simple nombre entier en un \emph{invariant topologique} : c'est le genre, augmenté de $1$, d'une surface de Riemann algébrique canoniquement associée à la cascade. Cette traduction porte trois conséquences importantes.

\begin{enumerate}[nosep,leftmargin=*]
\item Le triplet $\{3, 5, 7\}$ n'est pas simplement compatible avec l'équation $\mu = \sum_{p \in S} p$ : il est l'unique support d'une courbe algébrique non triviale dans la classe de Kontsevich--Norbury. La rigidité du choix des premiers actifs reçoit ainsi une justification \emph{topologique}, complémentaire de la justification arithmétique fournie par les axiomes $\BAthree$--$\BAfour$.
\item La cascade PT s'inscrit dans le cadre général de la récursion topologique d'Eynard--Orantin : les invariants $\omega_{g,n}^{\Cper}$ associés héritent des identités classiques (relations de Virasoro, formules d'intersection sur $\overline{\mathcal M}_{g,n}$). Une substitution structurelle Mirzakhani $\to$ persistance, $2\pi^{2}/3 \to (1/2) \sum_{p \in S} \gammap{p}$, transporte les volumes de Weil--Petersson vers les invariants de la PT~\cite[\S 8]{PT_GeoFlow_PersistenceCurve}.
\item Le genre $14$ est, à la connaissance de l'auteur, le premier invariant topologique strictement non trivial associé à la structure arithmétique de la PT, ouvrant la voie à une étude cohomologique de ses applications.
\end{enumerate}
La section suivante développe une seconde lecture, complémentaire, qui équipe l'espace des paramètres $\mu$ d'une métrique riemannienne canonique vérifiant une équation de soliton de Ricci.

\section{Lecture différentielle : le quasi-soliton de Ricci}
\label{sec:soliton}

La cascade arithmétique de la PT induit naturellement, sur l'espace des paramètres $\mu$, une métrique riemannienne canonique. Cette section établit que cette métrique satisfait, asymptotiquement, l'équation du soliton de Ricci de Hamilton et Perelman, avec un coefficient explicite et exact. La PT s'inscrit ainsi dans le grand cadre de la géométrie différentielle moderne, et son point fixe arithmétique $\mustar$ y reçoit une lecture géométrique complémentaire de la lecture algébrique de la \ref{sec:courbe}. Tous les résultats sont rigoureusement prouvés dans~\cite{PT_GeoFlow_RicciSoliton}.

\subsection{La métrique Fisher--Bianchi induite par la cascade}
\label{sec:soliton:metric}

Soit $\mu > 0$ le paramètre arithmétique et $q = q(\mu) = 1 - 2/\mu$ le paramètre de la branche-sommet. Pour chaque premier actif $p \in S^{*} := \{3, 5, 7\}$, on définit le \emph{taux de Hubble local}
\begin{equation}
H_p(\mu) \;=\; -\dfrac{1}{2\mu^{2}}\,\dfrac{p\,q^{p-1}}{1 - q^{p}},
\label{eq:Hubble_local}
\end{equation}
qui mesure le taux de variation de la phase d'holonomie au premier~$p$ avec $\mu$. La \emph{métrique de Fisher--Bianchi de la PT}, notée $\gPT$, est la métrique riemannienne sur la variété temporelle modélisée $\mathbb R \times \mathbb R^{|S^{*}|}$ définie, dans la jauge $N(\mu) = \mu$, par
\begin{equation}
\gPT \;=\; -d\mu^{2} \;+\; \sum_{p \in S^{*}} a_p(\mu)^{2}\, dx_p^{2},
\qquad
a_p(\mu) := \exp\!\Bigl(\!\int H_p(\mu)\, d\mu\Bigr).
\label{eq:gPT}
\end{equation}
Cette métrique est l'analogue PT d'une métrique de Bianchi~I lorentzienne~\cite{ChowKnopf2004}, mais ses facteurs d'échelle $a_p(\mu)$ ne sont \emph{pas isotropes} : ils diffèrent d'un premier actif à l'autre, ce qui distingue $\gPT$ d'une métrique de Robertson--Walker. C'est précisément cette anisotropie qui permettra de déterminer un résidu de Ricci unique.

L'introduction de la métrique~\eqref{eq:gPT} est canonique au sens suivant : ses facteurs d'échelle dérivent directement de la formule d'holonomie $\BAthree$, et la jauge $N(\mu) = \mu$ est la seule qui rende le potentiel d'entropie $S = -\ln \alpha_{\rm sieve}$ une fonction \emph{régulière} de la coordonnée temporelle (\cite[\S 2]{PT_GeoFlow_RicciSoliton}).

\subsection{L'équation du soliton de Ricci}
\label{sec:soliton:eq}

Soit $f$ une fonction lisse sur la variété temporelle. L'équation \emph{du soliton de Ricci gradient} de Hamilton et Perelman~\cite{Hamilton1982,Perelman2002} s'écrit
\begin{equation}
\Ric(g) + \Hess(f) \;=\; \lambda\, g, \qquad \lambda \in \mathbb R.
\label{eq:soliton_eq}
\end{equation}
Les solutions stationnaires ($\lambda = 0$, dites \emph{steady}), expansives ($\lambda < 0$) et contractantes ($\lambda > 0$) classifient les régimes asymptotiques de l'évolution sous le flot de Ricci $\partial_t g = -2\,\Ric(g)$. La géométrie de Perelman étudie systématiquement ces solutions et leurs singularités, en particulier dans la preuve de la conjecture de Poincaré~\cite{Perelman2002}.

Dans le cadre PT, on choisit $f := S = -\ln \alpha_{\rm sieve}$ : la fonction de potentiel est l'opposé du logarithme du couplage de Pontryagin. L'équation~\eqref{eq:soliton_eq}, projetée sur les composantes purement spatiales $(pp)$, $p \in S^{*}$, donne trois équations en deux inconnues $(\dot f, \lambda)$. Le théorème principal exploite cette surdétermination.

\subsection{Théorème du quasi-soliton expansif asymptotique}
\label{sec:soliton:thm}

\begin{theoreme}[Quasi-soliton de Ricci asymptotique]
\label{thm:soliton_PT}
Soit $\gPT$ la métrique de Fisher--Bianchi de la PT~\eqref{eq:gPT} dans la jauge $N(\mu) = \mu$, et soit $f = S = -\ln \alpha_{\rm sieve}(\mu)$. En imposant la compatibilité des trois équations $(pp)$ de l'équation soliton~\eqref{eq:soliton_eq} pour déterminer $\lambda(\mu)$, on a
\begin{equation}
\boxed{\;
\lambda(\mu) \;=\; -\dfrac{1}{\mu^{4}} \;+\; O\!\Bigl(\dfrac{1}{\mu^{5}}\Bigr)
\quad\text{lorsque } \mu \to \infty.
\;}
\label{eq:lambda_main}
\end{equation}
Le coefficient $-1$ du terme dominant en $\mu^{-4}$ est \emph{exact} et \emph{indépendant} du choix de paire $(p_1, p_2)$ de premiers actifs utilisée pour résoudre le système.
\end{theoreme}

\textbf{Conséquence immédiate.} $\lambda < 0$ asymptotiquement : la métrique de Fisher--Bianchi de la PT est un \emph{quasi-soliton expansif} au sens de Hamilton--Perelman. Le qualificatif \og quasi-\fg{} renvoie au fait que $\lambda(\mu) \to 0$ lorsque $\mu \to \infty$ : un soliton steady strict aurait $\lambda \equiv 0$, et un soliton expansif strict aurait $\lambda$ constant non nul ; la PT relève d'un régime intermédiaire dont l'écart à $\lambda = 0$ décroît en $\mu^{-4}$.

\subsection{Esquisse de preuve : équivalence anisotropie--unicité}
\label{sec:soliton:proof}

La détermination de $\lambda(\mu)$ s'appuie sur une équivalence structurelle entre anisotropie effective des taux de Hubble et unicité de la solution $\lambda$.

\begin{proposition}[Anisotropie \texorpdfstring{$\Leftrightarrow$}{<=>} unicité de $\lambda$]
\label{prop:anisotropy}
Pour la métrique $\gPT$ et tout $\mu > 0$, les deux propriétés suivantes sont équivalentes :
\begin{enumerate}[nosep,label=\textup{(\roman*)}]
\item Il existe $p_1 \neq p_2$ dans $S^{*}$ tels que $H_{p_1}(\mu) \neq H_{p_2}(\mu)$ (anisotropie effective).
\item Le système des équations $(pp)$ pour $p \in S^{*}$ admet une solution \emph{unique} en l'inconnue $\lambda$.
\end{enumerate}
\end{proposition}

\textbf{Démonstration.} Lorsque les $H_p$ sont tous égaux, le système des trois équations $(pp)$ se réduit à une seule équation à deux inconnues $(\dot f, \lambda)$ : $\lambda$ est alors un degré de liberté libre. Dès qu'apparaît une seule paire $(p_1, p_2)$ avec $H_{p_1} \neq H_{p_2}$, la différence des deux équations $(p_1 p_1) - (p_2 p_2)$ élimine $\dot f$ et fournit $\lambda$ explicitement~\cite[\S 4]{PT_GeoFlow_RicciSoliton}. $\square$

\begin{lemme}[Anisotropie asymptotique automatique]
\label{lem:H_asymptotic}
Pour $p_1 \neq p_2$ dans $S^{*}$,
\[
H_{p_1}(\mu) - H_{p_2}(\mu) \;=\; \dfrac{p_2 - p_1}{\mu^{3}} + O\!\bigl(\mu^{-4}\bigr) \;\neq\; 0
\quad\text{pour tout } \mu \text{ suffisamment grand.}
\]
\end{lemme}

Le \ref{thm:soliton_PT} découle alors directement : pour $\mu \to \infty$, l'anisotropie effective est automatique (lemme), donc $\lambda$ est unique (proposition), et le calcul du terme dominant de la différence des équations $(pp)$ donne $\lambda(\mu) = -1/\mu^{4} + O(\mu^{-5})$. L'indépendance du coefficient $-1$ vis-à-vis du choix de $(p_1, p_2)$ résulte de la structure du système : tout choix de paire produit la même équation linéaire après élimination de $\dot f$, à des termes d'ordre $\mu^{-5}$ près.

\subsection{Interprétation : résidu, dynamique et lien Perelman}
\label{sec:soliton:perelman}

\textbf{Le résidu comme signature de dynamique non-triviale.}
Un soliton steady au sens strict satisferait $\Ric + \Hess(f) = 0$, c'est-à-dire que la métrique serait invariante sous le flot de Ricci (à difféomorphisme et reparamétrisation près). Le résidu non nul $\lambda(\mu) \sim -1/\mu^{4}$ de la PT montre que la cascade arithmétique \emph{n'est pas} dans un tel régime d'équilibre exact : elle évolue lentement vers l'équilibre, à un taux qui décroît avec la profondeur du crible. C'est une signature géométrique de \emph{dynamique non triviale} ; on peut, par analogie cosmologique, lire le résidu comme un \og temps de relaxation \fg{} interne à la cascade.

\textbf{Inscription dans le programme Hamilton--Perelman.}
Les solitons de Ricci sont les objets centraux du programme de Hamilton sur le flot de Ricci, formalisé et complété par Perelman~\cite{Perelman2002}. Ils servent à la classification des singularités du flot et à la preuve des conjectures de Poincaré et de géométrisation de Thurston. Le \ref{thm:soliton_PT} place ainsi la PT dans ce paysage : la cascade arithmétique des nombres premiers actifs y apparaît non plus comme un simple comptage, mais comme une \emph{structure géométrique différentielle} satisfaisant une équation d'évolution canonique.

\textbf{Articulation avec la lecture algébrique.}
Les deux lectures de l'objet PT --- algébrique (\ref{sec:courbe}) et différentielle (cette section) --- ne se contentent pas de coexister : elles sont reliées par un \emph{pont d'échelle} démontré dans~\cite[\S 6]{PT_GeoFlow_RicciSoliton} :
\[
\lambda(\mu) \;\sim\; B_0(\mu)^{-8/3}, \qquad \mu \to \infty,
\]
où $B_0(\mu)$ est le coefficient dominant de l'invariant $\omega_{1,1}^{\Cper}$ de la récursion topologique de la courbe de persistance. L'exposant $-8/3$ est vérifié numériquement à $0{,}4\,\%$ ; son identité \emph{exacte} reste une question ouverte, qui constituera l'un des objectifs du programme spectral présenté en \ref{sec:programme}.

Avec ces deux lectures en place, la section suivante développe une troisième : la lecture spectrale, qui relie directement l'axiome arithmétique $\sper = 1/2$ aux signatures spectrales du couplage de Higgs et des zéros de Riemann.

\section{Lecture spectrale : identités K2 et A6}
\label{sec:spectrales}

Les deux lectures précédentes (algébrique en \ref{sec:courbe}, différentielle en \ref{sec:soliton}) montrent que la cascade arithmétique de la PT admet des réalisations géométriques canoniques. La présente section établit une troisième lecture : la cascade laisse des \emph{signatures spectrales} dans deux structures classiques de la mathématique pure et de la physique théorique --- la distribution des zéros de la fonction zêta de Riemann d'une part, et le couplage de Higgs du Modèle Standard d'autre part. Le résultat principal (identité K2) est une triple identification du résidu de Riemann--von Mangoldt comme expression de l'axiome arithmétique $\sper = 1/2$ ; le résultat auxiliaire (identité de trace A6) est une formule exacte vérifiée numériquement à $10^{-25}$. Ces résultats sont issus de la consolidation~\cite{PT_NewMath_Consolidation}.

\subsection{Densité des zéros de Riemann et résidu \texorpdfstring{$1/8$}{1/8}}
\label{sec:spectrales:RvM}

Soit $N(T)$ le nombre de zéros non triviaux $\rho = \beta + i\gamma$ de $\zeta(s)$ avec $0 < \gamma < T$. Le théorème classique de Riemann et von~Mangoldt~\cite{Edwards1974} fournit l'expansion asymptotique
\begin{equation}
N(T) \;=\; \dfrac{T}{2\pi}\,\log\!\dfrac{T}{2\pi e} \;+\; \dfrac{7}{8} \;+\; S(T) \;+\; O(1/T),
\label{eq:RvM}
\end{equation}
où $S(T) = \frac{1}{\pi}\,\arg \zeta(1/2 + iT)$ est la fonction d'argument et oscille de moyenne nulle. La constante $7/8$ est une signature spectrale fine~: elle code, en termes de la quantification semi-classique de Berry--Keating~\cite{BerryKeating1999}, une \emph{phase de Maslov} associée aux points tournants du modèle hamiltonien $H_{\rm BK} = xp$ régularisé.

La PT \emph{prédit}~\cite{PT_NewMath_Consolidation} une variante semi-classique de cette formule, où $7/8 = 1 - 1/8$ est remplacé par $1$ exactement~:
\begin{equation}
N_{\rm PT}(T) \;=\; \dfrac{T}{2\pi}\,\log\!\dfrac{T}{2\pi e} \;+\; 1 \;+\; S(T) \;+\; O(1/T).
\label{eq:NPT}
\end{equation}
L'écart entre~\eqref{eq:RvM} et~\eqref{eq:NPT} est exactement $\Delta_N = 7/8 - 1 = -1/8$, ou en valeur absolue, $|\Delta_N| = 1/8$. La question posée par l'identité K2 est : ce résidu $1/8$ a-t-il une lecture PT pure ?

\subsection{Identité K2 : triple identification du résidu \texorpdfstring{$1/8$}{1/8}}
\label{sec:spectrales:K2}

\begin{theoreme}[Identité K2, identification spectrale-physique]
\label{thm:K2}
Le résidu absolu $|\Delta_N| = 1/8$ admet, au point fixe $\mustar = 15$ de la cascade de la PT, trois identifications exactes structurellement équivalentes :
\begin{equation}
\boxed{\;
\dfrac{1}{8} \;=\; \dfrac{\sper^{2}}{2} \;=\; \lambdaH \;=\; \DeltaMaslov \;,
\;}
\label{eq:K2}
\end{equation}
où $\lambdaH$ est l'auto-couplage du boson de Higgs du Modèle Standard (résultat \texttt{R\,lambda\_H} de l'article compagnon~\cite{companion_M5}), et $\DeltaMaslov$ est la phase de Maslov des coins de la double barrière du modèle de Berry--Keating régularisé.
\end{theoreme}

\textbf{Trois identifications.} L'identité~\eqref{eq:K2} se déchiffre en trois sous-identités exactes, chacune dérivable indépendamment dans le cadre PT.
\begin{enumerate}[nosep,leftmargin=*,label=\textbf{P\arabic*.}]
\item \emph{Carré du paramètre de symétrie.} L'axiome $\sper = 1/2$ donne immédiatement
\[
\dfrac{\sper^{2}}{2} \;=\; \dfrac{1/4}{2} \;=\; \dfrac{1}{8}.
\]
\item \emph{Auto-couplage de Higgs.} La dérivation du couplage scalaire de Higgs développée dans l'article compagnon~\cite{companion_M5} fournit, à profondeur cascade $d = 3$,
\[
\lambdaH \;=\; \dfrac{\sper^{2}}{2} \;=\; \dfrac{1}{8}.
\]
Cette identité est obtenue indépendamment de toute considération spectrale, via la structure d'autocouplage du potentiel scalaire au point fixe.
\item \emph{Phase de Maslov.} Le modèle de Berry--Keating régularisé pour les zéros de Riemann donne, pour la phase de Maslov des coins de la double barrière~\cite{BerryKeating1999},
\[
\DeltaMaslov \;=\; \dfrac{1}{8} \;=\; -\dfrac{\pi/8}{\pi/1},
\]
en lecture semi-classique WKB, indépendamment de toute identification physique.
\end{enumerate}

\textbf{Origine commune.} Les trois identifications~P1--P3 partagent une origine structurelle commune : la bifurcation $\qplus / \qminus$ du paramètre de la loi géométrique au point fixe~$\mustar = 15$. Cette bifurcation produit, dans la PT, la même structure $\sper^{2}/2$ qui apparaît à la fois comme auto-couplage scalaire (côté physique du pont) et comme phase de Maslov (côté spectral). L'identité K2 affirme l'\emph{unicité} de cette origine~: les trois manifestations sont projections d'un même objet PT sur trois axes différents (arithmétique, physique, spectral).

\textbf{Statut.} L'identification K2 est démontrée \emph{structurellement} : chacune des trois sous-identités P1, P2, P3 est exacte, et leur coïncidence numérique au point fixe en fait une identité PT \emph{plausible} au sens du programme. La dérivation cohomologique (item K3 du programme~\cite{PT_NewMath_Consolidation}, en cours) fournirait une preuve unifiée des trois, à partir d'un même invariant cohomologique de l'involution $D : \qplus \leftrightarrow \qminus$.

\subsection{Identité de trace A6}
\label{sec:spectrales:A6}

\begin{theoreme}[Identité de trace exacte en $\Re(s) > 1$]
\label{thm:A6}
Soit $R(s)$ le facteur résiduel d'une représentation Fredholm de la fonction zêta régularisée par la PT, et soit $A_p$ l'amplitude de l'opérateur canalisé au premier $p$, $A_p = \sin^{2}\thetap{p}(\qplus) + \sin^{2}\thetap{p}(\qminus)$. Dans la région $\Re(s) > 1$ on a l'identité analytique exacte
\begin{equation}
\boxed{\;
-\dfrac{d}{ds} \log R(s) \;=\; \sum_{p} \log p \cdot \biggl[\, \dfrac{1}{p^{s} - 1} \,+\, A_p \cdot p^{-s} \,\biggr] \,. \;}
\label{eq:A6}
\end{equation}
\end{theoreme}

\textbf{Vérification numérique.} L'identité~\eqref{eq:A6} a été testée numériquement (multi-précision \texttt{mpmath}, $25$ chiffres significatifs) sur 17 valeurs de $s$ réparties dans $\Re(s) > 1$ ; l'écart entre les deux membres est au niveau de la précision machine, soit \emph{de l'ordre de} $10^{-25}$. À toutes fins pratiques, l'identité est \emph{exacte}~\cite[Note 62]{PT_NewMath_Consolidation}.

\textbf{Lecture.} Le premier terme $\sum_p \log p / (p^{s} - 1)$ est la dérivée logarithmique standard de la fonction zêta de Riemann, $-\zeta'(s)/\zeta(s)$. Le second terme, $\sum_p \log p \cdot A_p \cdot p^{-s}$, est la \emph{contribution PT propre} : il code l'effet de la bifurcation $\qplus / \qminus$ sur la trace régularisée. L'identité~\eqref{eq:A6} affirme donc que la dérivée logarithmique du facteur résiduel $R(s)$ se décompose proprement en un \og fond zêta classique \fg{} et une \og correction PT \fg{} d'amplitude $A_p$.

\subsection{Conséquence : lien direct \texorpdfstring{$\sper = 1/2$}{s=1/2} \texorpdfstring{$\leftrightarrow$}{<->} phénoménologie spectrale}
\label{sec:spectrales:conclusion}

Les deux identités K2 et A6 fournissent ensemble un \emph{pont arithmétique-spectral} concret~:
\begin{itemize}[nosep,leftmargin=*]
\item l'identité K2 (\ref{thm:K2}) relie une \emph{quantité observable} ($\lambdaH$, mesurable au LHC) à une \emph{quantité spectrale} ($\DeltaMaslov$, prédictible à partir des zéros de Riemann) en passant par une \emph{quantité PT pure} ($\sper^{2}/2$ = carré de l'unique entrée formelle de la théorie) ;
\item l'identité de trace A6 (\ref{thm:A6}) fournit la \emph{plomberie analytique} qui transporte les amplitudes PT $A_p$ dans la dérivée logarithmique d'une fonction résiduelle proche de $\zeta$, en régime $\Re(s) > 1$.
\end{itemize}
Ces deux résultats sont, au sens strict, des \emph{identités prouvées} (et non des théorèmes inconditionnels au sens de la \ref{sec:courbe} et de la \ref{sec:soliton}) : K2 repose sur la coïncidence structurelle de trois sous-identités, A6 est analytique en $\Re(s) > 1$ mais son prolongement à la ligne critique demeure ouvert. C'est précisément ce prolongement qui fait l'objet du programme spectral-arithmétique présenté en \ref{sec:programme}.

\section{Programme spectral-arithmétique : la conjecture RH-PT}
\label{sec:programme}

Les identités K2 (\ref{thm:K2}) et A6 (\ref{thm:A6}) sont les pierres d'angle d'un programme de recherche plus ambitieux, en cours de développement~\cite{PT_NewMath_Consolidation}, dont l'horizon est une approche PT-native de l'hypothèse de Riemann. Cette section esquisse le programme, en distinguant rigoureusement ce qui est démontré, ce qui est conjectural, et ce qui reste ouvert. \emph{Aucun} des résultats de cette section n'est revendiqué comme théorème inconditionnel.

\subsection{Opérateur de survie résiduel}
\label{sec:programme:operateur}

Soit, pour chaque premier $p$ et chaque complexe $s$ avec $\Re(s) > 1$, l'amplitude résiduelle
\[
\kappa_p(s) \;=\; 1 \;-\; \bigl(1 - p^{-s}\bigr)\,\exp\!\bigl(A_p\, p^{-s}\bigr),
\]
où $A_p = \sin^{2}\thetap{p}(\qplus) + \sin^{2}\thetap{p}(\qminus)$. La fonction zêta régularisée résiduelle de la PT s'écrit comme un produit eulérien modifié,
\[
\Rres(s) \;=\; \prod_{p} \dfrac{1}{1 - \kappa_p(s)},
\qquad \Re(s) > 1,
\]
qui converge absolument dans le demi-plan $\Re(s) > 1$. La dérivée logarithmique de $\Rres$ fait précisément apparaître l'identité de trace A6 (\ref{eq:A6}). On définit alors l'\emph{opérateur de survie résiduel} de la PT comme l'opérateur diagonal sur $\ell^{2}(\mathcal P)$ de symbole spectral $\lambda_p(s) = p^{-s}$ ; cet opérateur joue, dans le cadre PT, un rôle analogue à celui du hamiltonien semi-classique $\HBK = xp$ de Berry--Keating~\cite{BerryKeating1999}.

\subsection{Statut analytique et frontière critique}
\label{sec:programme:statut}

\begin{description}[nosep,leftmargin=*]
\item[Théorème.] L'identité de Schur--Fredholm associée à l'opérateur de survie résiduel est démontrée comme \emph{identité opératorielle exacte} sur l'espace de Hilbert $\ell^{2}(\mathcal P)$ dans le régime $\Re(s) > 1$~\cite{PT_NewMath_Consolidation}.
\item[Identité A6.] La dérivée logarithmique de $\Rres(s)$ admet la décomposition exacte en deux composantes ($-\zeta'/\zeta$ + correction PT $A_p p^{-s}$) en $\Re(s) > 1$, vérifiée numériquement à $10^{-25}$ (théorème~\ref{thm:A6}).
\item[Identité K2.] La phase de Maslov de la quantification semi-classique des zéros de Riemann coïncide avec $\sper^{2}/2 = 1/8$, triple identification spectrale-physique-arithmétique (théorème~\ref{thm:K2}).
\item[Conjectural.] Le prolongement analytique de l'identité~\eqref{eq:A6} à la ligne critique $\Re(s) = 1/2$ est ouvert. Les premières conditions de bord (régularisation de Hadamard, poids de Abel, lissage gaussien) ont été testées numériquement~\cite{PT_NewMath_Consolidation} et reproduisent la position des zéros de Riemann à toutes les précisions disponibles ; mais aucune preuve analytique rigoureuse sur la ligne critique n'est connue à ce jour.
\end{description}

\subsection{La conjecture RH-PT}
\label{sec:programme:conjecture}

\begin{conjecture}[RH-PT, conjecture du verrou spectral]
\label{conj:RH_PT}
La distribution distributionnelle des pôles de la trace régularisée
\[
\mathcal T(s) \;:=\; \mathrm{Tr}_{\rm reg}\!\bigl[\,(s - i\,\HBK)^{-1} \cdot D_R(s)\,\bigr]
\]
sur la ligne critique $\Re(s) = 1/2$ coïncide \emph{exactement} avec l'ensemble des zéros non triviaux de la fonction $\zeta$ de Riemann.
\end{conjecture}

La conjecture RH-PT est testée numériquement sur 17 zéros, sur 12 décades, en multi-précision~\cite[Notes 66--69]{PT_NewMath_Consolidation}. À toutes les précisions disponibles la coïncidence est vérifiée, mais le statut épistémique demeure \emph{conjectural} : la preuve rigoureuse sur la ligne critique reste à fournir.

\subsection{Ce que le programme apporterait}
\label{sec:programme:portee}

Si la conjecture RH-PT était démontrée, plusieurs résultats classiques et nouveaux en découleraient :
\begin{enumerate}[nosep,leftmargin=*]
\item l'hypothèse de Riemann elle-même, comme conséquence du fait que les pôles de $\mathcal T$ sur la ligne critique seraient en correspondance bijective avec les zéros de $\zeta$ ;
\item une explication structurelle de la position $\Re(s) = 1/2$ : la ligne critique apparaîtrait comme le \emph{lieu spectral} où l'opérateur PT $\HBK$ admet la régularisation canonique ;
\item un cadre opératoriel précis pour le programme d'Hilbert--Pólya, dans lequel le rôle de l'opérateur auto-adjoint est joué par l'opérateur de survie résiduel de la PT.
\end{enumerate}

\textbf{Mise en garde.} Le programme RH-PT relève à ce jour de la \emph{recherche ouverte}, et son inclusion dans cet article ne préjuge en rien de sa résolution. Les théorèmes des sections~\ref{sec:courbe} et~\ref{sec:soliton}, ainsi que les identités~K2 et~A6 de la section~\ref{sec:spectrales}, sont \emph{indépendants} de la conjecture RH-PT et restent valides en l'état. Le présent article propose RH-PT comme un \emph{horizon de programme}, non comme un résultat acquis.

\section{Conclusion et perspectives}
\label{sec:conclusion2}

\textbf{Synthèse des trois lectures.} La cascade arithmétique des nombres premiers actifs $\{3, 5, 7\}$ de la PT, fixée par le point fixe $\mustar = 15$, admet trois réalisations canoniques de natures complémentaires.
\begin{description}[nosep,leftmargin=*]
\item[Lecture algébro-géométrique (\ref{sec:courbe}).] La cascade est l'unique courbe algébrique dans la classe de Kontsevich--Norbury satisfaisant les conditions (A) d'holonomie polynomiale, (B) de seuil d'activation et (C) d'auto-cohérence. Son genre vaut $\genus(\Cper) = \mustar - 1 = 14$.
\item[Lecture différentielle (\ref{sec:soliton}).] La métrique de Fisher--Bianchi $\gPT$ induite par la cascade satisfait l'équation du soliton de Ricci asymptotique $\Ric + \Hess(f) = \lambda(\mu)\, \gPT$ avec $\lambda(\mu) = -1/\mu^{4} + O(1/\mu^{5})$, coefficient $-1$ exact.
\item[Lecture spectrale (\ref{sec:spectrales}).] L'axiome $\sper = 1/2$ se manifeste dans le résidu de Riemann--von Mangoldt via l'identité triple $1/8 = \sper^{2}/2 = \lambdaH = \DeltaMaslov$ (K2), et la dérivée logarithmique du facteur résiduel admet la décomposition exacte $-d/ds \log R(s) = \sum_p \log p\,[1/(p^{s} - 1) + A_p\,p^{-s}]$ vérifiée à $10^{-25}$ en $\Re(s) > 1$ (A6).
\end{description}

\textbf{Statuts épistémiques différenciés.} Les résultats de cet article se rangent en trois catégories qu'il convient de bien distinguer.
\begin{enumerate}[nosep,leftmargin=*]
\item \emph{Théorèmes prouvés inconditionnellement} : le théorème de classification et la formule du genre (\ref{thm:courbe_classification}, \ref{thm:genre}) ; le théorème du quasi-soliton de Ricci (\ref{thm:soliton_PT}).
\item \emph{Identités prouvées} : K2 (\ref{thm:K2}) est démontrée structurellement comme coïncidence exacte de trois sous-identités indépendantes ; A6 (\ref{thm:A6}) est une identité analytique exacte en $\Re(s) > 1$, vérifiée numériquement à toutes les précisions disponibles.
\item \emph{Programme conjectural} : la conjecture RH-PT (\ref{conj:RH_PT}) est un programme de recherche ouverte, en cours d'instruction~\cite{PT_NewMath_Consolidation}. Sa résolution éventuelle entraînerait l'hypothèse de Riemann.
\end{enumerate}

\textbf{Articulation avec l'article compagnon.} Cet article suppose, pour son intelligibilité complète, la lecture de~\cite{PT_Article1_Foundations}, où sont établis les sept théorèmes structurants $\Tzero$--$\Tsix$ et les trois axiomes de pont prouvés $\BAthree$--$\BAfive$. Les lectures géométrique et spectrale développées ici en sont des prolongements non triviaux : elles ne modifient pas le noyau arithmétique de la PT, mais en révèlent les facettes cohomologique, différentielle et spectrale. Une lecture en sens inverse (de cet article vers l'article 1) est également possible : le genre topologique 14 et l'équation du soliton fournissent des contraintes \emph{a posteriori} sur la cascade qui renforcent l'unicité du choix des premiers actifs $\{3, 5, 7\}$.

\textbf{Perspectives.} Trois directions de recherche immédiates se dégagent de cet article.
\begin{itemize}[nosep,leftmargin=*]
\item \emph{Calcul exhaustif des invariants $\omega_{g,n}^{\Cper}$.} La récursion topologique d'Eynard--Orantin permet, en principe, de calculer tous les invariants $\omega_{g,n}^{\Cper}$ de la courbe de persistance. Les premiers calculs (\cite[\S 8]{PT_GeoFlow_PersistenceCurve}) confirment la substitution structurelle $2\pi^{2}/3 \to (1/2) \sum \gammap{p}$ par rapport à la courbe d'Airy de Mirzakhani. Étendre ce calcul à $\omega_{2,1}^{\Cper}$ et au-delà permettrait de tester de nouvelles identités cascade.
\item \emph{Identité exacte du pont d'échelle $\lambda \leftrightarrow B_0$.} La relation $\lambda(\mu) \sim B_0(\mu)^{-8/3}$ est vérifiée numériquement à $0{,}4\,\%$. L'identité \emph{exacte} (à coefficient près) entre le résidu de Ricci $\lambda$ et le coefficient dominant $B_0$ de la récursion topologique constituerait un pont algébrico-différentiel fort.
\item \emph{Promotion cohomologique de K2.} La dérivation de l'identité K2 à partir d'un invariant cohomologique de l'involution $D : \qplus \leftrightarrow \qminus$ (item K3 du programme~\cite{PT_NewMath_Consolidation}) fournirait une preuve unifiée des trois sous-identités, et un cadre rigoureux pour leur extension à des dimensions différentes de $\mathbb R^{4}$.
\end{itemize}

\textbf{Liens avec les autres applications de la PT.} Les lectures ici présentées sont en dialogue actif avec plusieurs domaines de la phénoménologie PT, en particulier la cosmologie (la métrique de Fisher--Bianchi est l'analogue PT de Bianchi~I, cf.~\cite{companion_M3,companion_M5}), la physique des hautes énergies (le couplage $\lambdaH$ identifié en K2 est l'auto-couplage de Higgs mesurable au LHC, cf.~\cite{companion_M5}), et la théorie analytique des nombres (la conjecture RH-PT relie l'opérateur de survie PT aux zéros de Riemann). Aucune de ces applications n'est développée ici ; elles constituent l'horizon \emph{aval} du présent article.

\textbf{Statut éditorial.} Le couple d'articles~\cite{PT_Article1_Foundations}--présent constitue, à la connaissance de l'auteur, le matériel d'entrée le plus consolidé pour une revue par les pairs externe de la Théorie de la Persistance. Les théorèmes T0--T6, les axiomes de pont BA0--BA5, la dérivation de $\alpha_{\rm EM}$, la courbe de persistance, le quasi-soliton de Ricci et les identités spectrales K2 et A6 y sont énoncés avec leur statut épistémique précis. Les questions restant ouvertes (programme RH-PT, identité exacte $\lambda \leftrightarrow B_0$, promotion cohomologique de K2) sont explicitement marquées comme telles et constituent l'horizon de recherche immédiat du programme.


\appendix
\section{Glossaire géométrique-spectral}
\label{app:glossaire}

Le tableau~\ref{tab:glossaire2} récapitule les notations géométriques et spectrales propres à cet article et leur lecture PT.

\begin{table}[htbp]
\centering
\small
\begin{tabular}{@{}p{3.5cm}p{5cm}p{6cm}@{}}
\toprule
Notation & Définition & Lecture PT \\
\midrule
$\Cper$ & Courbe algébrique de persistance dans $\mathcal{KN}$ & Unique solution de ACA$_{\mathcal P}$ avec $|S| \geq 2$ \\
$\genus(\Cper)$ & Genre de $\Cper$ & $\mustar - \lfloor |S|/2 \rfloor = 14$ \\
$\omega_{g,n}^{\Cper}$ & Invariants d'Eynard--Orantin & Cascade des corrélateurs PT \\
$\gPT$ & Métrique de Fisher--Bianchi & Métrique riemannienne canonique sur l'espace des $\mu$ \\
$\Ric(g)$ & Tenseur de Ricci de $g$ & Courbure différentielle de $\gPT$ \\
$\Hess(f)$ & Hessien de $f$ & Hessien du potentiel $f = -\ln \alpha_{\rm sieve}$ \\
$\lambda(\mu)$ & Résidu de Ricci & $-1/\mu^{4} + O(\mu^{-5})$, signature dynamique \\
$H_p(\mu)$ & Taux de Hubble local & Taux de variation de $\thetap{p}$ avec $\mu$ \\
$N(T)$ & Comptage de zéros de Riemann & Loi de Riemann--von~Mangoldt \\
$\Delta_N^{\rm Maslov}$ & Phase de Maslov des coins & Résidu spectral $1/8$ de RvM \\
$\lambdaH$ & Auto-couplage de Higgs & Couplage scalaire à profondeur cascade $d=3$ \\
$\Rres(s)$ & Facteur résiduel zêta régularisé & $\prod_p 1/(1 - \kappa_p(s))$, converge en $\Re(s) > 1$ \\
$\kappa_p(s)$ & Amplitude de canal au premier $p$ & $1 - (1 - p^{-s})\exp(A_p p^{-s})$ \\
$A_p$ & Amplitude PT à $p$ & $\sin^{2}\thetap{p}(\qplus) + \sin^{2}\thetap{p}(\qminus)$ \\
$\HBK$ & Hamiltonien de Berry--Keating régularisé & Opérateur PT $xp$ régularisé \\
\bottomrule
\end{tabular}
\caption{Glossaire géométrique-spectral de l'article 2.}
\label{tab:glossaire2}
\end{table}

\section{Table récapitulative des résultats}
\label{app:resultats}

Le tableau~\ref{tab:resultats} synthétise les résultats principaux de cet article, classés par statut épistémique et avec leurs références.

\begin{table}[htbp]
\centering
\small
\begin{tabular}{@{}p{3cm}p{6cm}p{2cm}p{3.5cm}@{}}
\toprule
Étiquette & Énoncé court & Section & Statut \\
\midrule
\multicolumn{4}{c}{\textsf{\textbf{Théorèmes prouvés inconditionnellement}}} \\[2pt]
Th.~Classification (\ref{thm:courbe_classification}) & $S = \{3, 5, 7\}$ unique solution ACA$_{\mathcal P}$ avec $|S| \geq 2$ & \ref{sec:courbe} & \textsc{[Thm]} \\
Th.~Algébricité (\ref{thm:algebraicite_PT}) & $\Cper$ algébrique, $Y^{2} \in \mathbb Q(x)$ & \ref{sec:courbe} & \textsc{[Thm]} \\
Th.~Genre (\ref{thm:genre}) & $\genus(\Cper) = \mustar - \lfloor |S|/2 \rfloor = 14$ & \ref{sec:courbe} & \textsc{[Thm]} \\
Th.~Quasi-soliton (\ref{thm:soliton_PT}) & $\lambda(\mu) = -1/\mu^{4} + O(\mu^{-5})$, coef.\ $-1$ exact & \ref{sec:soliton} & \textsc{[Thm]} \\
Prop.~Anisotropie (\ref{prop:anisotropy}) & Anisotropie $\Leftrightarrow$ unicité de $\lambda$ & \ref{sec:soliton} & \textsc{[Thm]} \\
\midrule
\multicolumn{4}{c}{\textsf{\textbf{Identités prouvées}}} \\[2pt]
K2 (\ref{thm:K2}) & $1/8 = \sper^{2}/2 = \lambdaH = \DeltaMaslov$ & \ref{sec:spectrales} & \textsc{[Id]} structurelle \\
A6 (\ref{thm:A6}) & $-d/ds \log R = \sum_p \log p\,[1/(p^{s}{-}1) + A_p p^{-s}]$ & \ref{sec:spectrales} & \textsc{[Id]}, $\Re(s) > 1$ \\
\midrule
\multicolumn{4}{c}{\textsf{\textbf{Programme conjectural}}} \\[2pt]
RH-PT (\ref{conj:RH_PT}) & Pôles de $\mathcal T(s)$ sur $\Re(s) = 1/2$ $\leftrightarrow$ zéros de $\zeta$ & \ref{sec:programme} & \textsc{[Conj]} \\
$\lambda \leftrightarrow B_0$ & $\lambda(\mu) \sim B_0(\mu)^{-8/3}$ exact & \ref{sec:soliton:perelman} & \textsc{[Conj]} \\
K3 cohomologique & Dérivation cohomologique unifiée de K2 & \ref{sec:spectrales:K2} & \textsc{[Open]} \\
\bottomrule
\end{tabular}
\caption{Tableau récapitulatif des résultats de l'article. \textsc{[Thm]} = théorème inconditionnel, \textsc{[Id]} = identité prouvée, \textsc{[Conj]} = conjecture, \textsc{[Open]} = problème ouvert.}
\label{tab:resultats}
\end{table}


\printbibliography

\end{document}
